% Standalone manuscript-ready draft.  Copy the material between
% \begin{document} and \end{document} into the main Overleaf manuscript.
% Required packages in the main manuscript:
%   \usepackage{amsmath,amssymb,booktabs,threeparttable,graphicx,array}

\documentclass[10pt]{article}
\usepackage[margin=1in]{geometry}
\usepackage{amsmath,amssymb,booktabs,threeparttable,graphicx,array}
\usepackage{microtype}

\newcommand{\Mthree}{\ensuremath{\mathrm{M3}}}
\newcommand{\Qthree}{\ensuremath{\mathrm{Q3}}}
\newcommand{\Qfive}{\ensuremath{\mathrm{Q5}}}

\begin{document}

\section{Design and stress testing of the advection initializer}
\label{sec:advection-initializer}

\subsection{Motivation and estimand}

The directional Vecchia approximation places lagged conditioning blocks in a
corridor oriented by an advection vector
$\boldsymbol v=(v_{\mathrm{lat}},v_{\mathrm{lon}})^{\mathsf T}$.
Although both components of $\boldsymbol v$ are estimated continuously in the
subsequent seven-parameter likelihood optimization, the conditioning graph
must first be constructed from an initial direction.  Rebuilding that graph
inside every optimizer iteration would be expensive and would also change the
approximate likelihood being optimized.  We therefore estimate a fast
advection anchor, construct the directional corridor once, and keep the graph
fixed throughout the final fit.  All competing initializers share the same
starting values for the five nuisance parameters and the same optimizer
budget; only the initial advection anchor differs.

Our baseline, denoted \Mthree, is a \emph{masked-FFT empirical argmin}.  The
label ``M3'' is an experiment identifier rather than a third-order model.  The
two refinements, \Qthree\ and \Qfive, use the same empirical surface as \Mthree
and only interpolate its local minimum below the grid resolution.

\subsection{Masked-FFT cross-semivariogram}

Let $Z_t(\boldsymbol x)$ be the field at hour $t$ on a regular
$n_1\times n_2$ grid, and let $M_t(\boldsymbol x)$ equal one when that grid cell
is observed and zero otherwise.  No spatial binning is required because every
candidate displacement is an integer grid offset.  To attenuate time-varying
additive level shifts, we center each hourly image using only its observed
cells,
\begin{equation}
  \bar Z_t
  =\frac{\sum_{\boldsymbol x}M_t(\boldsymbol x)Z_t(\boldsymbol x)}
  {\sum_{\boldsymbol x}M_t(\boldsymbol x)},
  \qquad
  Y_t(\boldsymbol x)
  =M_t(\boldsymbol x)\{Z_t(\boldsymbol x)-\bar Z_t\}.
  \label{eq:hourly-demean}
\end{equation}
This removes first-order temporal nonstationarity in the mean, but it does not
assume away temporal changes in the covariance or missingness pattern.

For a one-hour lag $\tau=1$ and spatial offset $\boldsymbol h$, the pooled
pairwise cross-semivariogram is
\begin{equation}
  \widehat\gamma_1(\boldsymbol h)
  =\frac{1}{2N(\boldsymbol h)}
  \sum_{t=1}^{T-1}\sum_{\boldsymbol x}
  M_t(\boldsymbol x)M_{t+1}(\boldsymbol x+\boldsymbol h)
  \left[\widetilde Z_t(\boldsymbol x)
  -\widetilde Z_{t+1}(\boldsymbol x+\boldsymbol h)\right]^2,
  \label{eq:cross-semivariogram}
\end{equation}
where $\widetilde Z_t=Z_t-\bar Z_t$ on observed cells and
\begin{equation}
  N(\boldsymbol h)=\sum_{t=1}^{T-1}\sum_{\boldsymbol x}
  M_t(\boldsymbol x)M_{t+1}(\boldsymbol x+\boldsymbol h).
\end{equation}
In the eight-hour application, pooling uses all seven adjacent hourly
transitions.  The lag $\tau=1$ therefore means one hour, not one spatial grid
cell.  We use the one-hour lag as the primary initializer because it provides
the largest number of overlapping pairs, avoids compounding displacement over
multiple hours, and generally retains the strongest temporal dependence;
longer lags can be treated as a sensitivity analysis when one-hour dependence
is weak.

Expanding the square in Equation~\eqref{eq:cross-semivariogram} gives four
linear cross-correlations for each adjacent pair:
\begin{align}
 N_t(\boldsymbol h)
   &=\sum_{\boldsymbol x}M_t(\boldsymbol x)
      M_{t+1}(\boldsymbol x+\boldsymbol h), \\
 A_t(\boldsymbol h)
   &=\sum_{\boldsymbol x}Y_t(\boldsymbol x)^2
      M_{t+1}(\boldsymbol x+\boldsymbol h), \\
 B_t(\boldsymbol h)
   &=\sum_{\boldsymbol x}M_t(\boldsymbol x)
      Y_{t+1}(\boldsymbol x+\boldsymbol h)^2, \\
 C_t(\boldsymbol h)
   &=\sum_{\boldsymbol x}Y_t(\boldsymbol x)
      Y_{t+1}(\boldsymbol x+\boldsymbol h),
 \label{eq:four-correlations}
\end{align}
so that
\begin{equation}
 \widehat\gamma_1(\boldsymbol h)
 =\frac{\sum_{t=1}^{T-1}
 \{A_t(\boldsymbol h)+B_t(\boldsymbol h)-2C_t(\boldsymbol h)\}}
 {2\sum_{t=1}^{T-1}N_t(\boldsymbol h)}.
 \label{eq:masked-gamma}
\end{equation}
For arrays $U$ and $V$, all displacements are obtained simultaneously using
the correlation theorem,
\begin{equation}
 \operatorname{Corr}(U,V)(\boldsymbol h)
 =\mathcal F^{-1}
 \left\{\overline{\mathcal F(U)}\,\mathcal F(V)\right\}
 (\boldsymbol h).
 \label{eq:fft-correlation}
\end{equation}
Thus, the four quantities in Equation~\eqref{eq:four-correlations} are computed
by applying Equation~\eqref{eq:fft-correlation} to the data, squared data, and
observation masks.  Missing observations are set to zero only inside the FFT;
the mask correlation $N_t(\boldsymbol h)$ restores the exact number of valid
pairs at every displacement.  This is analogous in spirit to the observation-
window correction used by a debiased Whittle likelihood because both the data
and their observation window are transformed.  Here, however, the mask is used
for exact pairwise normalization of a cross-semivariogram, rather than for
correcting the expected periodogram.

Zero-padding is essential.  An unpadded discrete Fourier transform returns a
\emph{circular} correlation and would spuriously match the eastern and western
boundaries, or the northern and southern boundaries.  We instead compute the
full linear correlation by padding each dimension to at least
$(2n_1-1)\times(2n_2-1)$ (an FFT implementation may pad further to an efficient
transform length), and then extract the entries centered at
$(n_1-1,n_2-1)+\boldsymbol h$.  We searched integer offsets
$h_j\in\{-20,\ldots,20\}$, retained cells with at least 1,000 pooled valid
pairs, and applied a missing-aware Gaussian smoother with bandwidth
$0.063^\circ$.  For the GEMS grid spacings
$(\Delta_{\mathrm{lat}},\Delta_{\mathrm{lon}})=(0.044^\circ,0.063^\circ)$,
the search region is approximately
$[-0.88^\circ,0.88^\circ]\times[-1.26^\circ,1.26^\circ]$.

The \Mthree\ anchor is the discrete minimum
\begin{equation}
 \widehat{\boldsymbol v}_{\Mthree}
 =\underset{\boldsymbol h\in\mathcal H}{\arg\min}\;
 \widehat\gamma_1(\boldsymbol h).
 \label{eq:m3}
\end{equation}
The direct calculation costs
$O\{(T-1)|\mathcal H|n_1n_2\}$ operations, whereas the masked-FFT calculation
costs $O\{(T-1)n_1n_2\log(n_1n_2)\}$ and returns the entire displacement
surface at once.

\subsection{Safeguarded local quadratic refinements}

Because Equation~\eqref{eq:m3} is restricted to integer offsets, its error has
a grid-quantization component.  Let $(i_0,j_0)$ denote the \Mthree\ cell and
write local coordinates in \emph{grid-cell units},
$\boldsymbol d=(u,v)^{\mathsf T}$, rather than degrees.  Both refinements fit
\begin{equation}
 q(\boldsymbol d)
 =\beta_0+\boldsymbol g^{\mathsf T}\boldsymbol d
 +\frac{1}{2}\boldsymbol d^{\mathsf T}H\boldsymbol d,
 \qquad
 \widehat{\boldsymbol d}=-\widehat H^{-1}\widehat{\boldsymbol g},
 \label{eq:quadratic}
\end{equation}
where $H$ is a symmetric $2\times2$ Hessian.  There are six distinct
coefficients: an intercept, two slopes, and the three unique elements of $H$.
The eigenvalues of $H$ are the curvatures along the two principal directions
of the local surface.

The \Qthree\ method uses ordinary least squares on the $3\times3$ patch
$u,v\in\{-1,0,1\}$, giving nine surface values and three residual degrees of
freedom.  This is a deterministic local interpolation, not a regression model
for nine independent observations: neighboring semivariogram cells are
correlated and each cell already averages many data pairs.  The \Qfive\ method
uses the $5\times5$ patch $u,v\in\{-2,-1,0,1,2\}$ and weighted least squares
with prespecified Gaussian weights
\begin{equation}
 w(u,v)=\exp\left\{-\frac{u^2+v^2}{2s^2}\right\},
 \qquad s=1.25\ \text{grid cells},
 \label{eq:q5-weight}
\end{equation}
leaving 19 nominal residual degrees of freedom.  The weighting stabilizes the
fit while limiting contamination from nonquadratic behavior or a neighboring
basin farther from the discrete minimum.

The same numerical safeguards are imposed on \Qthree\ and \Qfive:
\begin{enumerate}
 \item the complete local patch must lie inside the searched surface and all
       patch values must be finite;
 \item $\widehat H$ must be positive definite, i.e., both eigenvalues must be
       strictly positive;
 \item the Hessian condition number, calculated in grid-cell coordinates,
       must satisfy $\kappa(\widehat H)\leq100$; and
 \item the interpolated correction must remain inside the central half cell,
       $|\widehat d_{\mathrm{lat}}|\leq0.5$ and
       $|\widehat d_{\mathrm{lon}}|\leq0.5$.
\end{enumerate}
If any condition fails, the method returns the unmodified \Mthree\ anchor.
The physical correction is
$(\Delta_{\mathrm{lat}}\widehat d_{\mathrm{lat}},
\Delta_{\mathrm{lon}}\widehat d_{\mathrm{lon}})$.
Scaling before the fit is important: otherwise the unequal latitude and
longitude spacings can distort the numerical condition number without changing
the underlying surface geometry.

We also examined a leave-one-transition-out stability diagnostic as a possible
hard gate between \Mthree\ and \Qthree.  The gate was not retained as a primary
method: among the 12 downstream comparisons on which it changed the selected
anchor, \Qthree\ and \Mthree\ each achieved the lower common NLL six times, and
the gate rejected \Qthree\ in the most consequential case in which it avoided
a poor local solution.  The diagnostic may still be reported as a continuous
measure of surface stability, but it is not used to override the numerical
safeguards above.

\subsection{Stress-test design}

The initializer-only factorial experiment crossed eight non-grid-aligned
directions $(22.5^\circ+45^\circ k)$, three speeds
$(0.75,2,4)$ grid cells per hour, three signal regimes, and five independent
replicates, producing $8\times3\times3\times5=360$ data sets.  Relative to the
reference covariance, the strong, reference, and weak regimes multiplied
$(\text{range}_{t},\text{nugget})$ by $(1.5,0.5)$, $(1,1)$, and $(0.5,3)$,
respectively.  This factorial stresses slow motion, large motion, weak temporal
dependence, high nugget, and all four quadrants.

A downstream experiment then used a partially completed set of 30 synthetic
data sets.  For each data set we used three common nuisance starts: the truth;
a short-time/high-nugget start with multipliers
$(0.75,1.25,1.25,0.50,1.50)$; and a long-time/low-nugget start with multipliers
$(1.25,0.75,0.75,1.50,0.50)$ for
$(\sigma^2,\rho_{\mathrm{lat}},\rho_{\mathrm{lon}},\rho_t,\tau^2)$.
This produced 90 seven-parameter fits per initializer.  The partial run covered
all selected $22.5^\circ$ cases and the 0.75- and 2-cell cases at
$112.5^\circ$, but not the opposite quadrants, and is therefore reported as
descriptive interim evidence.  Final estimates were evaluated on a common
oracle-direction corridor because raw likelihoods from different Vecchia
factorizations are not directly comparable.

We reran the complete 360-data-set initializer factorial using the final
cell-scaled safeguards for both \Qthree\ and \Qfive.  The more expensive
downstream full-fit stress test predated \Qfive; hence its 30-data-set results
compare \Mthree\ and \Qthree, followed by a deliberately small downstream
\Qfive\ screen.  This distinction prevents preliminary full-fit evidence for
\Qfive\ from being interpreted as confirmatory.

\begin{table}[t]
\centering
\caption{Stress tests of the empirical and sub-grid advection initializers.}
\label{tab:advection-stress}
\begin{threeparttable}
\small
\textbf{A. Initializer-only factorial stress test (360 synthetic data sets)}\\[2pt]
\begin{tabular}{lrrrrrr}
\toprule
Method & Mean error & Median error & Angle error & Correct & Q accepted & Time \\
       & (cells)    & (cells)      & (degrees)   & quadrant &            & (ms) \\
\midrule
\Mthree        & 0.587 & 0.559 & 13.88 & 0.656 & ---   & 2.828 \\
\Mthree+\Qthree & 0.381 & 0.266 &  8.83 & 0.944 & 0.953 & 2.913 \\
\Mthree+\Qfive  & 0.391 & 0.285 &  9.11 & 0.944 & 0.958 & 2.914 \\
\bottomrule
\end{tabular}

\vspace{6pt}
\textbf{B. Partial fixed-corridor full-fit stress test
(30 data sets, three nuisance starts, 90 fits per method)}\\[2pt]
\begin{tabular}{lrrrrrr}
\toprule
Method & Common & Seven-parameter & Final advection & Gradient & Fit time & Start-var. \\
       & regret & error & error & convergence & (s) & trace \\
\midrule
\Mthree         & 0.000739 & 0.1790 & 0.02184 & 0.444 & 7.135 & 0.3208 \\
\Mthree+\Qthree & 0.000500 & 0.1697 & 0.01366 & 0.467 & 7.152 & 0.1634 \\
\bottomrule
\end{tabular}

\vspace{6pt}
\textbf{C. Preliminary \Qfive\ screen (four synthetic data sets)}\\[2pt]
\begin{tabular}{lrrrrr}
\toprule
Method & Seed error & Common NLL & Seven-parameter & Final advection & Fit time \\
       & (cells)    &            & error           & error           & (s) \\
\midrule
\Mthree          & 0.680 & 1.258142 & 0.591 & 0.1812 & 4.086 \\
\Mthree+\Qthree  & 0.604 & 1.257321 & 0.366 & 0.0307 & 4.050 \\
\Mthree+\Qfive   & 0.530 & 1.256986 & 0.396 & 0.0214 & 4.042 \\
\bottomrule
\end{tabular}
\begin{tablenotes}[flushleft]\footnotesize
\item Errors in panel A are measured in grid-cell units.  ``Start-var. trace''
is the trace of the standardized final-parameter covariance across the three
nuisance starts.  Panel B is an unbalanced partial run; 87--90\% of fits reached
the 20-evaluation limit.  In paired dataset-cluster bootstraps, the
\Mthree+\Qthree\ minus \Mthree\ differences were $-0.000240$
($95\%$ CI $[-0.000880,0.000146]$) for common NLL, $-0.00929$
($[-0.03928,0.01018]$) for seven-parameter error, and $-0.00818$
($[-0.02502,0.00079]$) for final advection error.  All intervals include zero.
All panel C fits reached the 12-evaluation limit, so panel C compares equal
budgets rather than converged estimators.
\end{tablenotes}
\end{threeparttable}
\end{table}

\subsection{Real-data timing and interpretation}

On three full GEMS days ($114\times159=18{,}126$ grid cells per hour and eight
hours per day), the mean initializer times were 0.0199, 0.0201, and 0.0200
seconds for \Mthree, \Mthree+\Qthree, and \Mthree+\Qfive, respectively.
Both quadratic refinements fell back to \Mthree\ on one of the three days; when
accepted, the largest correction was 0.312 grid cells.  On the selected
non-fallback day, a short seven-parameter fit required 40.41, 40.49, and 40.72
seconds, respectively.  The common-corridor NLLs were 1.403918, 1.403625, and
1.403915, with \Mthree+\Qthree\ smallest.  All three real-data fits reached the
eight-evaluation limit, so these values are timing and stability checks rather
than evidence of optimizer convergence.

The 360-data-set initializer stress test shows that both safeguarded local
quadratic corrections remove substantial grid quantization error at negligible
computational cost.  In this larger screen, \Qthree\ had slightly smaller mean
and median errors than \Qfive; \Qfive\ was nevertheless competitive and had a
slightly higher safeguard acceptance rate.  The partial downstream experiment
also favors \Qthree\ in its point estimates and suggests reduced sensitivity to
nuisance starts, but its confidence intervals include zero and its directional
coverage is incomplete.  In the smaller downstream screen, \Qfive\ had the
smallest final advection error but not the smallest seven-parameter error.
Accordingly, we treat \Mthree\ as the primary baseline, safeguarded \Qthree\ as
the primary refinement, and weighted \Qfive\ as a prespecified competing
refinement for the final confirmatory experiment.

\end{document}
